\documentclass[../Main.tex]{subfiles}

\begin{document}
\section{Vector Space of Linear Maps}
\begin{definition}{Linear map space}
    Let $V$ and $W$ be vector spaces. The \underline{linear map space} $\L(V, W)$ is:
    \begin{equation*}
        \L(V, W) = \subsetselect{\alpha : V \mapsto W}{\alpha \text{ is linear}}
    \end{equation*}
\end{definition}
\begin{definition}{Matrix of a linear map}
    The \underline{matrix of a linear map} $\alpha : V \mapsto W$ with respect to the ordered bases $B$ of $V$ and $C$ of $W$ is:
    \begin{equation*}
        [\alpha]^B_C = \left(\left([\alpha(b_i)]_C\right)_j\right)_{ij}
    \end{equation*}
    that is, the $i, j$ element of $[\alpha]^B_C$ is the $j$th element of the coordinate vector of $\alpha(b_i)$.
\end{definition}
\begin{theorem}
    Let $V, W$ be finite-dimensional vector spaces with $\dim(V) = n$, $\dim(W) = m$. Let $\alpha : V \mapsto W$ be a linear map.
    \begin{enumerate}
        \item For all $v \in V$, $[\alpha(v)]_C = [\alpha]_C^B [v]_B$;
        \item $[\alpha]^B_C$ is the only matrix that satisfies part 1;
        \item There exists an isomorphism:
            \begin{align*}
                E^B_C : \L(V, W) &\mapsto M_{m \times n}(\F)\\
                \alpha &\mapsto [\alpha]^B_C
            \end{align*}
    \end{enumerate}
    \label{thmLinMapProps}
\end{theorem}
\begin{proof}
    Again argue in terms of bases.
\end{proof}
\section{Properties of Matrices}
\subsection{Change of Basis}
\begin{proposition}
    Let $U, V, W$ be finite-dimensional vector spaces with bases $B, C, D$ respectively. Let $\alpha \in \L(U, V)$ and $\beta \in \L(V, W)$. Then:
    \begin{equation*}
        [\beta \circ \alpha]^B_D = [\beta]^C_D [\alpha]^B_C
    \end{equation*}
    \label{propComposeMatrices}
\end{proposition}
\begin{proof}
    Use \thmref{thmLinMapProps}.
\end{proof}
\begin{definition}{Change of basis matrix}
    For a vector space $V$ with $n$ dimensions, the \underline{change of basis matrix} from a basis $B$ to $B'$ is given by:
    \begin{equation*}
        P = [\Id_V]^B_{B'}
    \end{equation*}
\end{definition}
\begin{propositions}{
        Let $V, B, B'$ be defined as above. Let $W$ be an $m$-dimensional vector space.
        \label{propsChangeBasisCompose}
    }
    \item $[\Id_V]^B_{B'}$ is invertible with inverse $[\Id_V]^{B'}_B$;
    \item if $\alpha \in \L(V, W)$ and $C, C'$ are bases for $W$ then:
        \begin{equation*}
            [\alpha]^{B'}_{C'} = [\Id_W]^{C}_{C'} [\alpha]^B_C [\Id_V]^{B'}_B
        \end{equation*}
\end{propositions}
\begin{proof}
    Use \thmref{propComposeMatrices}.
\end{proof}
\begin{definition}{Equivalence of matrices}
    Matrices $A, A' \in M_{m\times n}(\F)$ are \underline{equivalent} if there exist $P \in GL_m(\F)$ and $Q \in GL_n(\F)$ such that $A' = PAQ$.
\end{definition}
\begin{theorem}
    Let $V$ and $W$ be vector spaces with dimensions $n$ and $m$. Let $\alpha : V \mapsto W$ be a linear map with rank $r$. Then there exists bases $B, C$ for $V$ and $W$ such that:
    \begin{equation*}
        [\alpha]^B_C =
        \begin{pmatrix}
            I_r & 0 \\
            0 & 0
        \end{pmatrix}
    \end{equation*}
    and if for some bases $B', C'$ we find:
    \begin{equation*}
        [\alpha]^{B'}_{C'} =
        \begin{pmatrix}
            I_{r'} & 0 \\
            0 & 0
        \end{pmatrix}
    \end{equation*}
    then $r' = r$.
    \label{thmLinMapIdentityBasis}
\end{theorem}
\begin{proof}
    Argue by bases, and use \thmref{thmRankNullity}.
\end{proof}
\subsection{Rank of a Matrix}
\begin{definition}{Column space}
    For $A \in M_{m \times n}(\F)$, the \underline{column space} of $A$, $\Col(A)$ is the subspace of $\F^m$ spanned by the columns of $A$. The \underline{column rank} of $A$ is $\crk(A) = \dim(\Col(A))$.
\end{definition}
\begin{definition}{Row space}
    For $A \in M_{m \times n}(\F)$, the \underline{row space} of $A$, $\Row(A)$ is the subspace of $\F^n$ spanned by the rows of $A$. The \underline{row rank} of $A$ is $\rrk(A) = \dim(\Row(A))$.
\end{definition}
\begin{lemma}
    For $A \in M_{m \times n}(\F)$ and $B \in M_{n \times p}(\F)$, representing the linear maps $\alpha$ and $\beta$ in standard basis,
    \begin{equation*}
        \rk(\alpha \circ \beta) \leq \min\{\rk(\alpha), \rk(\beta)\}
    \end{equation*}
    \label{lemRankComposition}
\end{lemma}
\begin{proof}
    $\im(\alpha \circ \beta) \leq \im(\alpha)$ so we have the first condition. Next suppose $v \in \ker{\beta}$ and so is also in the kernel of $\alpha \circ \beta$. Use this logic and \thmref{thmRankNullity} for the required result.
\end{proof}
\begin{theorem}
    Let $A, A' \in M_{m \times n}(\F)$. Let $r = \crk(A)$.
    \begin{enumerate}
        \item $A$ is equivalent to:
            \begin{equation*}
                \begin{pmatrix}
                    I_r & 0 \\
                    0 & 0
                \end{pmatrix}
            \end{equation*}
        \item $A$ and $A'$ are equivalent if and only if $\crk(A) = \crk(A')$.
    \end{enumerate}
    \label{thmCrkEquivalences}
\end{theorem}
\begin{proof}
    Let $\alpha$ represent $A$ in standard basis. Then by \thmref{thmLinMapIdentityBasis} there exist bases for which $\alpha$ has the required form, so let $P$ and $Q$ be the required change of basis matrices for equivalence.

    Use the first part to show that if the ranks match, both are equivalent to the identity form in part 1. If instead we have equivalence, use \thmref{lemRankComposition} twice to bound the ranks above and below by each other.
\end{proof}
\begin{theorem}[Column rank is row rank]
    For any $A \in M_{m \times n}(\F)$,
    \begin{equation*}
        \rrk(A) = \crk(A)
    \end{equation*}
    and thus it makes sense to define $\rk(A)$.
    \label{thmCrkIsRrk}
\end{theorem}
\begin{proof}
    By \thmref{thmCrkEquivalences}, we find $P$, $Q$ such that:
    \begin{equation*}
        PAQ =
        \begin{pmatrix}
            I_{\crk(A)} & 0 \\
            0 & 0
        \end{pmatrix}
    \end{equation*}
    Taking the transpose does not affect the RHS, and we get $Q^TA^TP^T$. $Q$ and $P$ are still invertible when transposed, so $A^T$ is also equivalent to the matrix above.
\end{proof}
\subsection{Matrix Similarity}
\begin{definition}{Matrix similarity}
    Matrices $A, A' \in M_{n \times n}(\F)$ are \underline{similar} if there exists $P \in GL_n(\F)$ such that:
    \begin{equation*}
        A' = P^{-1} A Q
    \end{equation*}
\end{definition}
\section{Elementary Operations on Matrices}
The three types of \underline{elementary row operations}, with their corresponding\\\underline{elementary matrices} in brackets, are:
\begin{enumerate}
    \item swapping rows $i$ and $j$ ($T_{i,j}$),
    \item multiplying elements of row $i$ by $\lambda$ ($M_{i,\lambda}$),
    \item adding $\lambda$ times row $j$ to row $i$ ($C_{i, j, \lambda}$).
\end{enumerate}
The corresponding \underline{elementary column operations} are similar.

For elementary matrix $E$, apply $E$ to the left for a row operation and to the right for a column operation. $C_{i, j, \lambda}$ is different - it will add $\lambda$ times column $i$ to column $j$.

Elementary operations preserve the rank of a matrix.
\begin{definition}{Row-reduced echelon form}
    A matrix $A \in M_{m\times n}$ is in \underline{row-reduced echelon form} (RRE form) if:
    \begin{enumerate}
        \item All non-zero rows of $A$ appear above all zero rows of $A$.
        \item The leftmost non-zero element of a non-zero row is 1. This is the \underline{pivot entry} $1 \leq p(\vec{r_i}) \leq n$.
        \item If $\vec{r_i}, \vec{r_j} \neq \vec0$ and $i < j$ then
        \begin{equation*}
            p(\vec{r_i}) < p(\vec{r_j}).
        \end{equation*}
        \item In a column containing a pivot entry, every other entry is zero.
    \end{enumerate}
\end{definition}
A matrix is in CRE form if its transpose is in RRE form.
\begin{lemma}
    The rank of a matrix in RRE form is the number of non-zero rows it has.
    \label{lemRRERankIsNonZeroRows}
\end{lemma}
\begin{proof}
    Show that the non-zero rows form a basis of $\Row(A)$.
\end{proof}
\begin{theorem}
    Every matrix $A \in M_{n \times n}(\F)$ can be put into RRE form.
    \label{thmRRE}
\end{theorem}
\begin{proof}
    Let $\{\vec{c_1}, \cdots, \vec{c_{n}}\}$ be the columns of $A$.

    Induct on $n$: for the base case scale the first element to 1 and then subtract this from the other columns to get the rest zero.

    For the inductive step, first perform a transposition to get the first entry non-zero. Then scale this to be 1, and add this row to the rest to get all zeroes in the rest of the first column. Now apply the inductive hypothesis to the sub-matrix starting from row/column 2.
\end{proof}
\begin{theorem}
    For $A \in M_{n \times n}(\F)$, the following are equivalent:
    \begin{enumerate}
        \item $\rk(A) = n$,
        \item $A$ is a product of elementary matrices,
        \item $A$ is invertible.
    \end{enumerate}
    \label{thmSquareRREProps}
\end{theorem}
\begin{proof}
    First, if $A$ has full rank then \thmref{thmRRE} gives us that $A$ is the identity matrix in RRE form. Therefore, we can show that $A$ is a product of elementary matrices.

    Then, if $A$ is a product of elementary matrices we can invert the individual matrices to get that $A$ itself is invertible.

    Finally, if $A$ is invertible then consider the linear map that it represents. By \thmref{corLinPigeon} we have that it is surjective and so $A$ has rank $n$.
\end{proof}
\end{document}